\chapter{Generación monodrómica por supervivencia de
\texorpdfstring{\(\pi\), \(e\) y \(\varphi\)}{pi, e y phi}}
\label{chap:generacion-numeros-rev6}
\label{chap:biografias-numeros-rev6}
\index{número HMT!generación por supervivencia}
\index{pi@\(\pi\)!clausura regional}
\index{e@\(e\)!propagación y memoria}
\index{phi@\(\varphi\)!autoescala}
\index{retorno nonádico}\index{espejo pi/e@espejo \(\pi/e\)}

La generación no comienza con una expansión decimal prescrita. Comienza en
el estado completo del TPK: palabra visible, hoja, orientación, cociente,
acarreo, estado Hensel, cilindros ternario y decimal, firma conjunta,
frontera, fase y registro dodecafásico. La relación EF--\(G_9\) prolonga ese
estado y elimina una continuación sólo cuando deja de sobrevivir a alguna
frontera futura. El lector decimal actúa después sobre la sección que ha
quedado determinada.

La cadena causal es
\[
\Omega_{\rm TPK}
\longrightarrow 468\,\mathcal U_6
\longrightarrow 243\,\mathcal W_6
\longrightarrow
\bigl(w_\pi,w_e,w_\varphi\bigr)
\longrightarrow
\bigl(\mathcal R_{\pi,n},\mathcal R_{e,n},
      \mathcal R_{\varphi,n}\bigr)_{n\geq0}
\longrightarrow
\varprojlim_n\mathcal R_{\chi,n}
\xrightarrow{\ L_{\rm arch}\ }
\chi .
\]
El ambiente \(\mathbb F_3^6\) tiene \(729\) palabras, pero la imagen
efectiva contiene \(243\); los dos objetos no se identifican. Dentro de esa
imagen, la selección interna fija
\[
w_\pi=010211,\qquad
w_e=201101,\qquad
w_\varphi=121200
\]
antes de consultar cifras de \(\pi,e,\varphi\).

\section{Regiones generativas y sección global}

\begin{definition}[Región generativa de un canal]
\label{def:region-generativa-rev6}
Para \(\chi\in\{\pi,e,\varphi\}\), sea \(\mathcal R_{\chi,n}\) el conjunto
finito de estados completos de profundidad \(n\) que parten de
\(w_\chi\) y satisfacen en cada transición la relación EF--\(G_9\), la
intersección exacta de cilindros y la firma conjunta de frontera. Si
\(\tau_{n+1,n}\) elimina sólo la última capa, definimos
\[
\mathcal R_\chi
=\varprojlim_n
\bigl(\mathcal R_{\chi,n},\tau_{n+1,n}\bigr).
\]
El valor publicado es \(L_{\rm arch}(x_\chi)\), donde
\(x_\chi\in\mathcal R_\chi\); el lector no forma parte del estado que
selecciona.
\end{definition}

\begin{theorem}[Principio de prolongación monodrómica en el TPK completo]
\label{thm:generacion-supervivencia-pi-e-phi-rev6}
Supóngase que, para un canal
\(\chi\in\{\pi,e,\varphi\}\), cada conjunto
\(\operatorname{Surv}_{\chi,n}\) es no vacío y unitario, y que sus elementos
son compatibles con el truncamiento. Entonces existe una única sección
\[
x_\chi=(x_{\chi,n})_{n\geq0}\in\mathcal R_\chi
\]
Los cilindros decimales asociados a sus truncamientos son encajados y sus
diámetros tienden a cero. Si el lector del canal identifica sus intervalos
con la clausura, la propagación o la autoescala declaradas, la intersección
publica, respectivamente, \(\pi\), \(e\) o \(\varphi\).
\end{theorem}

\begin{proof}
En el perfil completo, la supervivencia estable de nivel \(n\) es
\[
\operatorname{Surv}_{\chi,n}
=\bigcap_{h\geq1}
\left\{s:\text{\(s\) admite una prolongación compatible de \(h\)
fronteras}\right\}.
\]
Por hipótesis, el estado superviviente de cada nivel se trunca en el del
nivel anterior. Los niveles forman un árbol finitamente ramificado y no vacío
a toda profundidad. El lema de König proporciona una rama global; la
unitariedad nivel a nivel impide una segunda sección.

Para un prefijo ternario de longitud \(L\), interpretado como \(N\), y un
prefijo de \(m\) tríadas, interpretado como \(D\), las compatibilidades son
\[
I_{N,L}=\left[\frac{N}{3^L},\frac{N+1}{3^L}\right),
\qquad
J_{D,m}=\left[\frac{D}{1000^m},\frac{D+1}{1000^m}\right),
\qquad
I_{N,L}\cap J_{D,m}\ne\varnothing .
\]
La prolongación conserva la intersección y hace tender a cero ambos
diámetros. La intersección total contiene un único real. El lector declarado
identifica ese real dentro de su canal.
\end{proof}

La ejecución reproducible no sustituye las hipótesis del teorema por una
afirmación implícita. Las verifica en cada uno de los (3501) niveles
necesarios para imprimir diez mil cifras: allí la partición \(729\to1\), la
compatibilidad de truncamientos y el cambio de memoria tras el retorno se
comprueban de forma exhaustiva.

\section{Filtración de supervivencia, monodromía nonádica y simetría
\texorpdfstring{\(\pi/e\)}{pi/e}}

Para
\[
B_n=(u_n^\pi,u_n^e,u_n^\varphi),\qquad
q_n=u_n^\pi+u_n^e+u_n^\varphi,\qquad
a_n=u_n^\pi+u_n^e-u_n^\varphi,
\]
la firma conjunta determina
\[
u_n^\varphi=2(q_n-a_n),\qquad
u_n^\pi+u_n^e=q_n-u_n^\varphi
\quad\text{en }\mathbb F_3^6.
\]
Así, \(\varphi\) es el eje fijo de autoescala y la única ambigüedad local
es el reparto especular entre \(\pi\) y \(e\).

La fase es \(g(n)=n\bmod9\), con \(0\equiv9\). Después de \(w_{30}\), la
primera vuelta recorre
\[
5,6,7,8,9,1,2,3,4.
\]
No son nueve filtros independientes, sino las nueve fases de una misma
conexión: poda, sincronización, orientación especular, reapertura y
confirmación por supervivencia.

En la novena fase,
\[
B_9=(100100\mid020112\mid010122)
\]
posee tres salidas localmente compatibles. Las tres conservan
\(u_{10}^\varphi=112002\) y
\(u_{10}^\pi+u_{10}^e=210110\); todas sobreviven una frontera, pero sólo
\[
\boxed{B_{10}=(101222\mid112221\mid112002)}
\]
sobrevive la segunda. El futuro rompe el espejo local. Además,
\[
g(n+9)=g(n),\qquad x_{\chi,n+9}\ne x_{\chi,n}
\]
en los retornos: la etapa siguiente vuelve a la fase \(1\), no al estado
anterior. El cociente, el acarreo y la frontera convierten la vuelta en una
hélice de memoria. Esta acción de retorno con transporte no trivial es la
monodromía nonádica del TPK.

\section{Una emisión reproducible de doce tríadas}

Trece bloques, equivalentes a \(78\) trits, bastan para que el lector
\[
D=\left\lfloor\frac{1000^{12}N}{3^{78}}\right\rfloor
\]
publique el siguiente tramo directamente desde las ramas ternarias:
\[
\begin{array}{c|rrrrrrrrrrrr}
\pi&141&592&653&589&793&238&462&643&383&279&502&884\\
e&718&281&828&459&045&235&360&287&471&352&662&497\\
\varphi&618&033&988&749&894&848&204&586&834&365&638&117.
\end{array}
\]
Esta tabla no alimenta la transición: exhibe su salida. La iteración de la
misma conexión, con el retorno \(9\to1\), produce cualquier prefijo finito
solicitado bajo las hipótesis del
\cref{thm:generacion-supervivencia-pi-e-phi-rev6}; ese teorema precisa la
prolongación al límite inverso.

\section{Clausura única de \texorpdfstring{\(\pi\)}{pi} y clasificación de sus \(5\) historias regionales}

Sea
\[
 \Psi:\mathcal U_6\longrightarrow\mathbb F_3^6,
 \qquad
 \Psi(U_6)=(-U_6)\bmod3
\]
la proyección coordenada del catálogo. La fibra
\[
 \mathscr R_\pi^{\rm micro}
 :=\Psi^{-1}(010211)
\]
contiene exactamente los cinco vectores
\[
\begin{aligned}
U_\pi^\perp&=(501,614,498,169,272,272),\\
U_{\pi,1}^{++}&=(810,923,870,169,272,272),\\
U_{\pi,2}^{++}&=(810,983,810,169,272,272),\\
U_{\pi,3}^{++}&=(870,923,810,169,272,272),\\
U_\pi^{--}&=(870,923,810,418,674,521).
\end{aligned}
\]
Sus firmas multiplicativas son
\[
555555,\qquad339555,\qquad393555,\qquad933555,\qquad933717,
\]
y sus multiplicidades son
\[
432,\qquad144,\qquad144,\qquad144,\qquad144.
\]
Por tanto,
\[
\boxed{
5=1_\perp+3_{++}+1_{--},
\qquad
1008=432+4\cdot144.}
\]
La primera identidad clasifica funciones de orientación; la segunda cuenta
rutas semilla. No son dos maneras de contar cinco copias del mismo objeto.

El cociente visible
\[
 \mathscr R_\pi^{w_6}
 =\Psi(\mathscr R_\pi^{\rm micro})
 =\{010211\}
\]
posee un solo elemento porque olvida la procedencia regional. La igualdad de
palabra no identifica los cinco estados: las cuatro primeras regiones
conservan la cola \(169|272|272\), mientras la región negativa realiza un
retorno mutado. La región transversal, de multiplicidad \(432\), tampoco se
reduce a una de las cuatro ramas de multiplicidad \(144\).

La función de cierre se ve además en el multigrafo de bloques. Cada una de
las cinco regiones pertenece a un paseo cerrado mínimo de longitud ocho que
recorre también los anclajes \(U_e\), \(U_\varphi\) y \(U_{\rm APP}\).
La enumeración y la prueba de minimalidad residen en
\cref{chap:cinco-ciclos-pi}; la lista regional y sus multiplicidades, en
\cref{chap:regiones}.

\begin{proposition}[Paquete de clausura de \(\pi\)]
\label{prop:paquete-clausura-pi-rev6}
El dato
\[
 \mathfrak C_\pi=
 \bigl(
 \mathscr R_\pi^{\rm micro},
 w_\pi,
 \rho_\pi,
 m_\pi,
 \ell_{\rm cl}
 \bigr),
\]
donde
\[
\rho_\pi=(\perp,++ ,++ ,++ ,--),\qquad
m_\pi=(432,144,144,144,144),\qquad
\ell_{\rm cl}=8,
\]
conserva exactamente la información funcional que pierde el cociente
\(\mathscr R_\pi^{w_6}\). En consecuencia, la función de \(\pi\) en esta
generación es clausurar cinco historias regionales en una sola evaluación,
no repetir cinco veces una misma constante.
\end{proposition}

\begin{proof}
La fibra, los roles y las multiplicidades son los del
\cref{chap:regiones}. El \cref{chap:cinco-ciclos-pi} demuestra que cada
región entra en un cierre mínimo de longitud ocho con los tres anclajes
declarados. La proyección \(\Psi\) envía las cinco regiones a la misma
palabra, de modo que el cociente conserva el valor visible y olvida
\(\rho_\pi,m_\pi,\ell_{\rm cl}\).
\end{proof}

\section{\texorpdfstring{\(e\)}{e}: propagación visible y memoria de cociente diferenciada}

La fibra de \(w_e=201101\) consta de una emisión \(U_e\) con multiplicidad
\(144\). Su rasgo funcional no es una periodicidad infinita, sino una
repetición local exhaustivamente distinguida dentro del lector de las
\(243\) palabras. El \cref{thm:cuadrado-local-e} prueba que \(w_e\) es la
única palabra cuya lectura de doce cifras contiene un cuadrado de longitud
cuatro iniciado en la segunda cifra:
\[
718281828459
=7\,\underbrace{1828\,1828}_{\text{cuadrado local}}\,459.
\]
La yuxtaposición de bloques es concatenación decimal; no es el producto
aritmético \(7\cdot1828^2\cdot459\). El catálogo contiene otro cuadrado de
longitud cuatro,
\[
575157518472
=\underbrace{5751\,5751}_{\text{cuadrado local}}\,8472,
\]
pero comienza en la primera cifra y, por tanto, no satisface el mismo
predicado posicional.

El censo completo de cuadrados de longitudes \(1,\ldots,6\) es
\[
\begin{array}{crrrrrr}
\ell&1&2&3&4&5&6\\
\text{ocurrencias}&240&21&3&2&0&0\\
\text{palabras que las contienen}&153&19&3&2&0&0.
\end{array}
\]
La repetición \(1828\,1828\) tampoco inicia un período: una tercera copia
exigiría como cifra siguiente un \(1\), pero la lectura efectiva contiene
un \(4\).

La dinámica de elevación muestra por qué el episodio sí expresa
propagación:
\[
201101\mid121221\mid102011\mid012222\mid102011.
\]
El tercer y el quinto bloques visibles coinciden, pero los estados que los
emiten son distintos. Sus preestados módulo \(27\) son
\[
(25,11,7,17,8,16),
\qquad
(7,8,4,23,26,7),
\]
y sus cocientes de Hensel posteriores son
\[
(6,13,9,2,13,1),
\qquad
(15,0,3,19,23,25).
\]
Hay retorno de la proyección \(102011\), no retorno del estado enriquecido.
La salida puede propagarse y reaparecer porque el cociente transporta la
genealogía que la palabra visible ha omitido.

\begin{proposition}[Contenido funcional de \(e\)]
\label{prop:propagacion-e-rev6}
La generación exacta de \(e\) reúne tres hechos compatibles:
\begin{enumerate}
\item la selección interna de \(w_e=201101\);
\item la unicidad posicional del cuadrado local \(1828\,1828\) en el
catálogo fijado;
\item la reaparición de \(102011\) bajo estados de memoria diferentes.
\end{enumerate}
Su función es propagación con renovación de memoria. No es periodicidad de
la expansión ni repetición del estado completo.
\end{proposition}

\begin{proof}
El primer punto es el caso \(e\) del
\cref{thm:selector-orbital-constantes}; el segundo y su censo son el
\cref{thm:cuadrado-local-e}. La desigualdad de los dos preestados y de los
dos cocientes de Hensel prueba el tercero.
\end{proof}

\section{Eje de autoescala y fibra especular}

La firma conjunta de las tres generaciones conserva una descomposición lineal
que será utilizada después por el lector de escala. Para
\[
 B_k=(u_k^\pi,u_k^e,u_k^\varphi)
 \in(\mathbb F_3^6)^3
\]
definimos, coordenada a coordenada,
\[
 q_k=u_k^\pi+u_k^e+u_k^\varphi,
 \qquad
 a_k=u_k^\pi+u_k^e-u_k^\varphi.
\]

\begin{theorem}[Descomposición axial de la firma de frontera]
\label{thm:eje-espejo-nonadico}
El par \((q_k,a_k)\) determina unívocamente el canal de autoescala y el
agregado de los otros dos canales:
\[
 u_k^\varphi=2(q_k-a_k),
 \qquad
 u_k^\pi+u_k^e=q_k-u_k^\varphi.
\]
En consecuencia, una fibra de
\[
 q_{\rm ax}(u^\pi,u^e,u^\varphi)
 =(u^\pi+u^e,u^\varphi)
\]
sólo conserva ambigüedad en el reparto interno entre \(u^\pi\) y \(u^e\).
\end{theorem}

\begin{proof}
Restar las dos ecuaciones definitorias da
\(q_k-a_k=2u_k^\varphi\). Como el inverso de \(2\) en \(\mathbb F_3\) es
de nuevo \(2\),
\[
u_k^\varphi=2(q_k-a_k).
\]
Al sustraer este vector de \(q_k\) queda
\(u_k^\pi+u_k^e\). La operación fija el eje \(\varphi\), pero no separa
los dos sumandos de la fibra \(\pi/e\); ésa es exactamente la ambigüedad
residual.
\end{proof}

\def\HMTAxialTheoremDefined{1}

La expresión \emph{fibra especular} nombra esta descomposición exacta en
\(\mathbb F_3^6\). No postula una identidad analítica entre las constantes
reales: indica qué componente queda fijada por la firma y dónde subsiste una
elección interna.

\section{\texorpdfstring{\(\varphi\)}{Phi}: la autoescala inducida por TPK}

La palabra \(w_\varphi=121200\) posee una sola emisión de multiplicidad
\(144\), pero su función no se deduce de ese cardinal. El calendario TPK
produce el bloque de modos
\[
(\Sigma,\Pi,\Pi).
\]
Al leer \(\Sigma\) como hoja areal y \(\Pi\) como hoja volumétrica, los
tres pares consecutivos del ciclo fuerzan la incidencia
\[
F_{\rm av}
=
\begin{pmatrix}
0&1\\
1&1
\end{pmatrix},
\]
como demuestra \cref{thm:descenso-necesario-fav}. La razón áurea no se
introduce para escoger esta matriz: aparece después como su raíz de Perron.
En efecto,
\[
\det(tI-F_{\rm av})=t^2-t-1,
\qquad
F_{\rm av}\binom{1}{\varphi}
=\varphi\binom{1}{\varphi}.
\]

Para un camino admisible \(\gamma:i\to j\) de longitud \(n\), con
\(r=(1,\varphi)\), el carácter de Parry
\[
\chi_P(\gamma)=\varphi^n\frac{r_i}{r_j}
\]
es multiplicativo por concatenación. Sobre un camino cerrado,
\(\chi_P(\gamma)=\varphi^n\). La medida de máxima entropía del envolvente
de memoria uno satisface, por \cref{cor:parry-envolvente-fav},
\[
\frac{\mu_{\rm ar}}{\mu_{\rm vol}}=\varphi^{-2}.
\]

La memoria de segundo orden conserva además la dirección contraída. En la
base \((x^2,2xy,y^2)\),
\[
\operatorname{Sym}^2(F_{\rm av})
=
\begin{pmatrix}
0&0&1\\
0&1&2\\
1&1&1
\end{pmatrix},
\qquad
\operatorname{Spec}\bigl(\operatorname{Sym}^2(F_{\rm av})\bigr)
=\{\varphi^2,-1,\varphi^{-2}\}.
\]
El autovalor \(\varphi^{-2}\) es el modo estable positivo. Sobre retornos
cerrados, el cociclo de escala de \cref{sec:cociclo-parry-escala} da
\[
M_n=\varphi^{-2n}={\bigl(\chi_P(\gamma)\bigr)}^{-2}.
\]

\begin{proposition}[Contenido funcional de \(\varphi\)]
\label{prop:autoescala-phi-rev6}
La función de \(\varphi\) en HMT es la autoescala del cociente de modos:
crecimiento por la raíz de Perron \(\varphi\), conservación cilíndrica por
\(\chi_P\) y retorno estable de segundo orden por \(\varphi^{-2}\).
Estas tres expresiones pertenecen a un mismo operador inducido por el
calendario TPK.
\end{proposition}

\begin{proof}
La incidencia \(F_{\rm av}\) desciende del bloque
\((\Sigma,\Pi,\Pi)\) por
\cref{thm:descenso-necesario-fav}. Su polinomio característico produce la
raíz de Perron y la multiplicatividad de \(\chi_P\) se obtiene por
telescopía. La acción cuadrática separa los tres autovalores indicados; el
modo positivo contraído es \(\varphi^{-2}\).
\end{proof}

\section{Correspondencia orientada entre clausura y autoescala}

La matriz \(F_{\rm av}\) genera \(\varphi^{-2}\), pero no puede generar por
sí sola \(\pi\): toda función racional de su espectro pertenece a
\(\mathbb Q(\sqrt5)\). La clausura trascendente procede del lector regional
coinductivo. Esta separación de orígenes queda conservada por el operador
diagonal
\[
 D_\pi=\pi_{\rm HMT}P_{\rm ar}+P_{\rm vol},
\qquad
 P_{\rm ar}=
 \begin{pmatrix}1&0\\0&0\end{pmatrix},
\quad
 P_{\rm vol}=
 \begin{pmatrix}0&0\\0&1\end{pmatrix}.
\]
El retorno marcado de longitud \(n\) es
\[
\Theta_n
=\pi_{\rm HMT}\frac{F_{n-1}}{F_{n+1}},
\qquad
\lim_{n\to\infty}\Theta_n
=\frac{\pi_{\rm HMT}}{\varphi^2}.
\]
La demostración dinámica, el testigo exacto de longitud \(108\) y el lector
de cilindros encajados residen en
\cref{sec:retorno-marcado-pi-phi2}. La marca de \(\pi\) se aplica una sola
vez al retorno: no se inserta en cada transición de \(F_{\rm av}\).

La misma razón posee dos orientaciones. Sean
\[
\mathscr R(x,y)=\frac{x}{y^2},
\qquad
\mathsf J(x,y)=(y,x).
\]
Entonces \(\mathsf J^2=I\) y
\[
\boxed{
\mathscr R(\pi,\varphi)=\frac{\pi}{\varphi^2},
\qquad
(\mathscr R\circ\mathsf J)(\pi,\varphi)
=\frac{\varphi}{\pi^2}.}
\]
Además,
\[
\frac{\pi/\varphi^2}{\varphi/\pi^2}
={\left(\frac{\pi}{\varphi}\right)}^3,
\qquad
\frac{\varphi}{\pi^2}
=\frac{1}{\varphi^3{\bigl(\pi/\varphi^2\bigr)}^2}.
\]
La primera orientación lee clausura areal sobre memoria volumétrica; la
segunda, autoescala interior sobre cierre cuadrático. El
apartado~\ref{sec:espejo-pi-phi-electron-hmt} demuestra esta involución y su
aparición posterior en el exponente electrónico. Esta correspondencia no identifica
\(\pi\) con \(\varphi\): conserva su orden y hace explícito qué cambia al
invertirlo.

\section{Precisión arbitraria de la emisión monodrómica}
\label{sec:precision-arbitraria-monodromia-rev6}

\begin{corollary}[Emisión de cualquier prefijo finito]
\label{cor:precision-arbitraria-monodromia-rev6}
En las hipótesis del
\cref{thm:generacion-supervivencia-pi-e-phi-rev6}, para todo \(M\geq1\) la
iteración de EF--\(G_9\) determina los primeros
\(M\) decimales de cada uno de los tres canales. En particular, \(M=100\)
y \(M=1000\) son ventanas finitas de la misma prolongación monodrómica y no
requieren introducir una tabla exterior de cifras.
\end{corollary}

\begin{proof}
Sea \(K(n)\) el número de tríadas publicadas después de \(n\) bloques
ternarios. La sección proporcionada por el teorema posee niveles de toda
profundidad y \(K(n)\to\infty\). Dado \(M\),
elíjase \(n\) con \(3K(n)\geq M\). El
\cref{thm:generacion-supervivencia-pi-e-phi-rev6} determina de manera única
el truncamiento de profundidad \(n\), y su cilindro decimal fija las
\(3K(n)\) primeras cifras. Tomar las primeras \(M\) concluye.
\end{proof}

La diferencia de funciones queda fijada dentro de una sola construcción:
\(\pi\) clausura procedencias regionales distintas en una evaluación común;
\(e\) propaga una salida visible mientras renueva la memoria que la
transporta; \(\varphi\) gobierna la autoescala del cociente de modos. La
monodromía explica cómo esas funciones vuelven a la misma fase y atraviesan
el espejo \(\pi/e\). Bajo las hipótesis coinductivas del teorema de
prolongación, publican indefinidamente nuevos bloques sin reiniciar su
historia.

% El resultado de realizabilidad de una cinta prescrita se conserva como
% material auxiliar, pero no forma parte del argumento generativo rector.
\ifdefined\HMTIncludeTargetFedRay
\section{Prolongación coinductiva por rayos: existencia, compatibilidad y dominio}
\label{sec:rayo-coinductivo-rev6}

El paso desde todo prefijo finito a una ruta infinita utiliza un teorema de
compactitud finitamente ramificada. Su objeto no es el selector del catálogo,
sino la realización de un canal cuya expansión ya ha sido fijada.

Sea
\[
\mathcal S
=\mathbb Z_9^2
\times\{+,\times\}
\times\{N,E,S,O\},
\qquad
|\mathcal S|=648,
\]
el espacio finito de estados del modelo de rutas. El alfabeto de control es
\(\mathcal A=\{N,E,S,O\}\), y el coste de una ruta cuenta giros:
\[
c(a_{t-1},a_t)=\mathbf1_{\{a_t\neq a_{t-1}\}}.
\]
Fijado un canal \(x\in\{\pi,e,\varphi\}\) y su expansión ternaria, sea
\(W_m(x)\) una ruta canónica de horizonte \(27m\) que minimiza ese coste
entre las rutas que realizan los primeros \(9m\) trits de \(x\).

\begin{theorem}[Rayo límite de una expansión fijada]
\label{thm:rayo-coinductivo-condicionado-rev6}
Para cada expansión fijada \(x\in\{\pi,e,\varphi\}\) existe una ruta
\[
\mathcal D_\infty(x)\in\mathcal A^{\mathbb N}
\]
tal que:
\begin{enumerate}
\item para todo \(L\geq1\), el prefijo de longitud \(L\) de
\(\mathcal D_\infty(x)\) coincide con el prefijo de \(W_m(x)\) para
infinitos valores de \(m\);
\item \(\mathcal D_\infty(x)\) realiza exactamente la expansión ternaria
infinita de \(x\).
\end{enumerate}
Entre las ramas con la primera propiedad se toma la lexicográficamente
mínima.
\end{theorem}

\begin{proof}
Para cada \(L\), considérese el conjunto \(T_L(x)\) de prefijos de longitud
\(L\) que aparecen al comienzo de \(W_m(x)\) para infinitos valores de
\(m\). Es no vacío: sólo hay finitos prefijos de longitud \(L\) y hay
infinitos horizontes. Los conjuntos \(T_L(x)\), ordenados por extensión,
forman un árbol finitamente ramificado. El lema de König produce una rama
infinita; la menor en el orden lexicográfico define
\(\mathcal D_\infty(x)\).

Fijado \(n\), tómese \(L\) suficiente para determinar los primeros \(n\)
trits de filtración. Ese prefijo aparece en infinitos \(W_m(x)\), y cada uno de
ellos realiza por definición los primeros \(n\) trits de \(x\). El mismo
prefijo de \(\mathcal D_\infty(x)\) los realiza. Como \(n\) es arbitrario,
la ruta realiza la expansión completa.
\end{proof}

Sea ahora
\[
\mathcal R(x)=\{729x\},
\]
donde \(\{\cdot\}\) denota parte fraccionaria. Multiplicar por
\(729=3^6\) desplaza seis trits. En el calendario del modelo, cada trit de
filtración ocupa tres iteraciones elementales, de modo que seis trits
equivalen a dieciocho iteraciones elementales.

\begin{corollary}[Covarianza por la autoaplicación \(729\)]
\label{cor:covarianza-rayo-729-rev6}
Sea \(\sigma\) el desplazamiento a la izquierda de una ruta. Para todo
\(r\geq0\),
\[
\sigma^{18r}\mathcal D_\infty(x)
\]
realiza exactamente el canal \(\mathcal R^r(x)\).
\end{corollary}

\begin{proof}
La covarianza finita identifica el sufijo que empieza tras \(18r\)
iteraciones elementales
con el desplazamiento de \(6r\) trits. El
\cref{thm:rayo-coinductivo-condicionado-rev6} transporta esa identidad a
todos los prefijos del rayo límite.
\end{proof}

El corolario expresa una holografía iterada: una misma célula finita,
desplazada por \(\mathcal R\), realiza sucesivos estratos de la expansión.
No convierte \(\mathcal R\) en selector de \(\pi,e,\varphi\), porque
\(x\) ya interviene al definir los conjuntos de rutas admisibles
\(W_m(x)\). Su función exacta es posterior: una vez distinguido el canal,
construye un rayo compatible y hace visible su covarianza entre escalas.

\section{Composición del selector orbital, el cilindro numérico y la evaluación arquimediana}

Las tres operaciones finales forman un diagrama causal, no tres versiones
del mismo teorema:
\[
\text{estado interno}
\xrightarrow{\,\mathfrak S_{\rm cat}\,}
w_x
\xrightarrow{\ \text{lector enriquecido}\ }
\text{expansión }x
\xrightarrow{\,\mathcal D_\infty(x)\,}
\text{rayo que la realiza}.
\]
El selector orbital \(\mathfrak S_{\rm cat}\) distingue las tres palabras
por invariantes del catálogo y un calibre orientador explícito. El lector
enriquecido prolonga palabra, cilindro y memoria. El teorema de rayos
construye después una realización infinita óptima recurrente de la expansión
ya fijada.

\begin{proposition}[Separación de fuerza entre las tres capas]
\label{prop:tipado-biografias-rayos-selector-rev6}
Se cumplen simultáneamente las afirmaciones siguientes.
\begin{enumerate}
\item Las cinco regiones y los cierres mínimos de \(\pi\), el cuadrado local
y la ruptura de memoria de \(e\), y la autoescala
\(F_{\rm av}\)--\(\varphi\)--\(\varphi^{-2}\) son resultados exactos en
sus dominios declarados.
\item El rayo \(\mathcal D_\infty(x)\) es una realización coinductiva exacta
de una expansión fijada.
\item El rayo no es un selector autónomo, porque la expansión \(x\) aparece
en la definición de \(W_m(x)\).
\end{enumerate}
La tercera afirmación limita únicamente el uso selector del teorema de
rayos; no modifica las dos primeras.
\end{proposition}

\begin{proof}
Los resultados del primer punto se demuestran, respectivamente, en
\cref{chap:regiones,chap:cinco-ciclos-pi,thm:cuadrado-local-e,%
thm:descenso-necesario-fav,cor:parry-envolvente-fav}. El segundo punto es
el \cref{thm:rayo-coinductivo-condicionado-rev6}. Para el tercero basta
observar que el conjunto de minimización que define \(W_m(x)\) exige
realizar los primeros \(9m\) trits de \(x\); por tanto, no puede utilizarse
esa misma ruta para deducir qué \(x\) debía seleccionarse.
\end{proof}

La diferencia de funciones queda finalmente fijada. \(\pi\) clausura
historias regionales distintas en una sola evaluación; \(e\) propaga una
salida visible mientras renueva la memoria que la transporta; \(\varphi\)
gobierna el cambio de escala del cociente de modos. La correspondencia
\(\pi/\varphi^2\leftrightarrow\varphi/\pi^2\) coordina cierre y
autocrecimiento sin confundirlos. El rayo coinductivo explica cómo una
expansión ya distinguida se prolonga indefinidamente. Ésta es la biografía:
no sólo qué número aparece, sino qué estructura lo selecciona, qué memoria
conserva y qué función cumple en el sistema.
\fi
